\documentclass[12pt]{article}
\usepackage[utf8]{inputenc}
\usepackage[T1]{fontenc}
\usepackage{hyperref}
\title{Alagwa Language Resources for K-12 Curriculum Development}
\author{Compiled by Claude Sonnet 5 (claude-sonnet-5) for Andr\'as Kornai}
\date{}
\begin{document}
\maketitle

\section{Introduction}
Alagwa (ISO 639-3: wbj; Glottocode alag1248) is a West-Rift Southern Cushitic language of the Kondoa district in the Dodoma region of Tanzania, around Kolo. Speaker estimates vary: Wikipedia gives about 53,000 (2009), Joshua Project 69,000, and a 2016 monograph about 10,000. Some speakers have shifted to Sandawe. Other names include Wasi, Waasi, Chasi and Kialagwa.

\section{Primary Educational Resources}
No primers, textbooks or school programmes for Alagwa turned up in the searches.

\section{Linguistic and Reference Materials}
\begin{itemize}
\item Wikipedia entry on the Alagwa language: \url{https://en.wikipedia.org/wiki/Alagwa_language}
\item Maarten Mous (2016), \emph{Alagwa: a South Cushitic language of Tanzania}, the main descriptive work (catalogue entry): \url{https://www.koeppe.de/titel_alagwa-a-south-cushitic-language-of-tanzania}
\item WALS language entry: \url{https://wals.info/languoid/lect/wals_code_agw}
\end{itemize}

\section{Religious and Community Resources}
\begin{itemize}
\item Joshua Project profile for the Alagwa people, who are reported as overwhelmingly Muslim: \url{https://joshuaproject.net/people_groups/10242}
\item Joshua Project language page for wbj: \url{https://joshuaproject.net/index.php/languages/wbj}
\end{itemize}

\section{Digital and Audio Resources}
\begin{itemize}
\item Global Recordings Network language page for wbj: \url{https://globalrecordings.net/en/language/wbj}
\end{itemize}

\section{Curriculum Development Context}
The Mous monograph is the main basis for any orthography or teaching grammar. The searches found no school materials. Language shift among the young makes documentation and teaching materials more urgent.
\end{document}